\scanpage{3}

% archive sheet number: 2

D'une loi de \uline{coalgèbre} \inserted{\uline{unitaire}} sur $A$, \deleted{\unsure{de|sur} $U$,} donnée par

\[
\begin{aligned}
\Delta_U &: U\unsure{\longrightarrow|-}U\times_A U\\
\unsure{J|I|\eta}_U &: U\longrightarrow A\qquad \text{augmentation (\unsure{surjective|surjection})}
\end{aligned}
\]

\deleted{compatibles en le sens qu'on \unsure{aura} si \unsure{une|d'une} augmentation}

\[
\deleted{\delta_U x=\Delta_U x-1}
\]

qui soit \uline{commutative}, \uline{associative}, d'une \uline{structure de $A$-algèbre} sur $U$ qui soit \uline{associative}, \uline{unitaire} (pas nécessairement commutative) et telle que $U\otimes_A U\to U$ soit un hom de coalgèbres unitaires [ou, ce qui revient au \unsure{même}, $\Delta_U$ est un hom d'algèbres associatives unitaires]. Enfin on suppose qu'il existe un hom \inserted{de \uline{coalgèbre} \uline{unitaire}} $x\to\check{x}$, ayant la propriété suivante :

\[
x\in U,\quad \Delta_U x\unsure{=|>}\sum y_i\otimes z_{\deleted{\unsure{j}}i}
\quad\Longrightarrow\quad \sum\deleted{\unsure{z_i\check{y}_i}}y_i\check{z}_i=\sum\deleted{\unsure{\check{z}_i y_i}}\check{y}_i z_i=J_U(x).1_U.
\]

Alors $x\to\check{x}$ est un antiautomorphisme involutif de la structure d'algèbre de $U$, et \deleted{\unsure{cela}} il est défini de façon unique par les conditions énoncées.

Soit maintenant $C$ une $A$-algèbre. On veut expliquer \unsure{que|ce que} signifie que $C$ admet $B$ comme ``groupe d'opérateurs à \deleted{droite} gauche'' dans la catégorie $\unsure{\mathcal{C}|\mathscr{C}}$. Cela signifie ceci :
